\documentclass[11pt]{amsart}

\usepackage[T1]{fontenc}
\usepackage{amsmath,amssymb,amsthm}
\usepackage{aliascnt}
\usepackage{microtype}
\usepackage[hidelinks]{hyperref}
\hypersetup{
  pdftitle={A cubic logarithmic upper bound for commutators close to the identity},
  pdfauthor={Hyra},
  pdfsubject={Quantitative commutators in operator algebras},
  pdfkeywords={commutator, bounded operator, Cuntz isometries, supersolution}
}
\allowdisplaybreaks
\numberwithin{equation}{section}

\newtheorem{theorem}{Theorem}[section]

\newaliascnt{proposition}{theorem}
\newtheorem{proposition}[proposition]{Proposition}
\aliascntresetthe{proposition}

\newaliascnt{lemma}{theorem}
\newtheorem{lemma}[lemma]{Lemma}
\aliascntresetthe{lemma}

\newaliascnt{corollary}{theorem}
\newtheorem{corollary}[corollary]{Corollary}
\aliascntresetthe{corollary}

\theoremstyle{remark}
\newaliascnt{remark}{theorem}
\newtheorem{remark}[remark]{Remark}
\aliascntresetthe{remark}

\newcommand{\BH}{\mathcal B(H)}
\newcommand{\one}{\mathbf 1}
\newcommand{\norm}[1]{\left\lVert #1\right\rVert}
\newcommand{\abs}[1]{\left\lvert #1\right\rvert}
\newcommand{\ind}[1]{\mathbf 1_{\{#1\}}}
\newcommand{\diag}{\operatorname{diag}}
\newcommand{\id}{\operatorname{id}}
\newcommand{\e}{\mathrm e}

\title[A cubic logarithmic upper bound]{A cubic logarithmic upper bound for commutators close to the identity}
\author{Hyra}
\thanks{Hyra is an AI research-agent system.  The proof originated in Hyra
search runs; the accompanying repository contribution records provenance and
the exact scope of the Lean formalization.}
\date{}

\subjclass[2020]{Primary 47A63; Secondary 46L05, 47A20}
\keywords{commutator, bounded operator, Cuntz isometries, nonlinear fixed point, discrete supersolution}

\begin{document}

\begin{abstract}
Let $H$ be an infinite-dimensional complex Hilbert space and, for
$0<\varepsilon<1$, define
\[
 m_H(\varepsilon)
 :=\inf\bigl\{\norm D\norm X:
 D,X\in\mathcal B(H),\ \norm{[D,X]-\one}\leq\varepsilon\bigr\}.
\]
Popa's lower bound is
$m_H(\varepsilon)\geq \frac12\log(1/\varepsilon)$, while Tao proved the
upper bound $m_H(\varepsilon)=O(\log^5(1/\varepsilon))$.  We retain Tao's
operator-matrix template and replace the worst-case estimate for its linear
right inverse by an entrywise supersolution argument adapted to the endpoint
forcing in the nonlinear system.  The key input is a sharp norm estimate for
matrices of the form $\bigl(\begin{smallmatrix}a&b\\c&a\end{smallmatrix}\bigr)$
and a profile whose square has a controlled discrete concavity defect.  This
produces solutions of size $O(\delta n)$ with $\delta\asymp n^{-3}$ and yields
\[
 m_H(\varepsilon)
 \leq 5.3\times10^6\left(\log\frac1\varepsilon\right)^3,
 \qquad 0<\varepsilon\leq\frac12.
\]
The proof is uniform in $\varepsilon$ and applies to arbitrary, not
necessarily separable, infinite-dimensional Hilbert spaces.
\end{abstract}

\maketitle

\section{Introduction}

Throughout, $H$ is an infinite-dimensional complex Hilbert space,
$\mathcal B(H)$ is the algebra of bounded operators on $H$ with its operator
norm, $\one$ denotes the identity, and
\[
 [D,X]:=DX-XD.
\]
All logarithms without a subscript are natural.

The Wintner--Wielandt theorem states that the identity of a unital Banach
algebra cannot be a commutator of two bounded elements
\cite{Wintner,Wielandt}.  For $\mathcal B(H)$, Brown and Pearcy showed that the
identity can nevertheless be approximated arbitrarily well in norm by bounded
commutators \cite{BrownPearcy}.  The quantitative problem is to determine the
least possible growth of the product of the two operator norms.  For
$0<\varepsilon<1$, set
\begin{equation}\label{eq:def-m}
 m_H(\varepsilon)
 :=\inf\left\{\norm D\norm X:
 D,X\in\mathcal B(H),\ \norm{[D,X]-\one}\leq\varepsilon\right\}.
\end{equation}
Popa proved the lower bound
\begin{equation}\label{eq:popa}
 m_H(\varepsilon)\geq\frac12\log\frac1\varepsilon;
\end{equation}
see \cite{Popa}, and \cite{Tao} for a short proof.  Tao refined the
Brown--Pearcy construction and obtained
\begin{equation}\label{eq:tao}
 m_H(\varepsilon)
 =O\!\left(\log^5\frac1\varepsilon\right),
 \qquad 0<\varepsilon\leq\frac12,
\end{equation}
see \cite{Tao}.  Krishna and Johnson extended the same exponent to a class of
unital $C^*$-algebras containing $\mathcal B(H)$ and the Cuntz algebra
$\mathcal O_2$ \cite{KrishnaJohnson}.  A publicly available 2026 manuscript, accompanied by a Lean formalization,
gives an $O(\log^4(1/\varepsilon))$ refinement
\cite{Bilich}.  The argument below gives the following cubic logarithmic
upper bound.

\begin{theorem}\label{thm:main}
Let $H$ be an infinite-dimensional complex Hilbert space.  For every
$0<\varepsilon\leq\frac12$ there exist $D_\varepsilon,X_\varepsilon\in
\mathcal B(H)$ such that
\begin{equation}\label{eq:main-error}
 \norm{[D_\varepsilon,X_\varepsilon]-\one}\leq\varepsilon
\end{equation}
and
\begin{equation}\label{eq:main-norm}
 \norm{D_\varepsilon}\norm{X_\varepsilon}
 \leq 5.3\times10^6
 \left(\log\frac1\varepsilon\right)^3.
\end{equation}
Consequently,
\[
 m_H(\varepsilon)
 =O\!\left(\log^3\frac1\varepsilon\right)
 \qquad(\varepsilon\downarrow0).
\]
\end{theorem}

The algebraic template is Tao's.  For a matrix size $n$ and a small parameter
$\delta$, one constructs $D,X\in M_n(\mathcal B(H))$ so that
$[D,X]-\one$ is supported in the last column.  A nonlinear system is solved to
annihilate rows $2,\ldots,n$ of that column, leaving only the $(1,n)$ corner.
A diagonal similarity then multiplies this corner by $2^{-(n-1)}$, while the
norm product is $O(\delta^{-1})$.

The exponent is therefore governed by the largest admissible value of
$\delta$.  The linear part of the nonlinear system has a right inverse
$R=L(I-E)^{-1}$ whose global operator norm is $O(n^2)$.  A direct use of this
worst-case norm loses two powers of $n$.  The actual forcing, however, is an
endpoint spike.  We dominate the Neumann series entrywise by a scalar profile
$\psi$ satisfying
\begin{equation}\label{eq:intro-super}
 s_i+\frac12\norm{
 \begin{pmatrix}
 \psi_i&\psi_{i-1}\\
 \psi_{i+1}&\psi_i
 \end{pmatrix}}
 \leq\psi_i.
\end{equation}
The sharp two-by-two estimate in Lemma~\ref{lem:two-by-two} reduces
\eqref{eq:intro-super} to a discrete concavity inequality for
$\Psi_i:=\psi_i^2$.  The quadratic profile
\begin{equation}\label{eq:intro-profile}
 \Psi_i=144\delta^2 i(2n-i)
\end{equation}
has a constant interior concavity defect and a much larger defect at the right
endpoint, exactly where the forcing is concentrated.  This gives
$\norm{\beta_i}=O(\delta n)$, permits $\delta\asymp n^{-3}$, and hence yields
the exponent three after taking $n\asymp\log(1/\varepsilon)$.

The proof establishes an upper bound only.  In particular, it does not show
that the exponent three is optimal; the lower and upper logarithmic exponents
remain one and three, respectively.

An accompanying Lean 4 project machine-checks the sharp two-by-two norm
estimate, the concavity and supersolution core, the quadratic-profile
arithmetic, and the final numerical constants.  It does not formalize the
operator-valued Neumann-series convergence, the Banach fixed-point
construction, or the complete existence theorem in $\mathcal B(H)$; those
steps remain part of the proof below.

\section{Operator-matrix preliminaries}

For every finite $n$, an infinite-dimensional Hilbert space is unitarily
isomorphic to its $n$-fold Hilbert direct sum.  Thus there is a unitary
$V_n:H^{\oplus n}\to H$, and conjugation by $V_n$ identifies
$M_n(\mathcal B(H))=\mathcal B(H^{\oplus n})$ isometrically and
$*$-isomorphically with $\mathcal B(H)$.  We shall construct the operators in
$M_n(\mathcal B(H))$ and use this identification at the end.

We record the elementary norm estimates used below.

\begin{lemma}\label{lem:matrix-norms}
Let $M=(M_{ij})\in M_n(\mathcal B(H))$.
\begin{enumerate}
\item If $M$ has only one nonzero entry, at $(p,q)$, then
$\norm M=\norm{M_{pq}}$.
\item If $M_{ij}=0$ for $j\neq q$, then
\[
 \norm M\leq\sum_{i=1}^n\norm{M_{iq}}.
\]
\item If, for a fixed integer $k$ with $1-n\leq k\leq n-1$,
$M_{ij}=c_i\ind{j=i+k}$, then
\[
 \norm M
 =\max_{\substack{1\leq i\leq n\\1\leq i+k\leq n}}
 \norm{c_i}.
\]
\item For $z_{11},z_{12},z_{21},z_{22}\in\mathcal B(H)$,
\[
 \norm{
 \begin{pmatrix}
 z_{11}&z_{12}\\z_{21}&z_{22}
 \end{pmatrix}}_{M_2(\mathcal B(H))}
 \leq
 \norm{
 \begin{pmatrix}
 \norm{z_{11}}&\norm{z_{12}}\\
 \norm{z_{21}}&\norm{z_{22}}
 \end{pmatrix}}_{M_2(\mathbb R)}.
\]
\item The operator norm of a real $2\times2$ matrix with nonnegative entries
is nondecreasing in each entry.
\end{enumerate}
\end{lemma}

\begin{proof}
The first statement follows by testing on unit vectors supported in the
$q$-th coordinate and using an approximating sequence for the norm of
$M_{pq}$.  The second and third statements follow directly from the action of
$M$ on $H^{\oplus n}$; in the third statement one tests on a single valid
coordinate and an approximating sequence of unit vectors in $H$ to obtain the
reverse inequality.

For the fourth statement, if $h=(h_1,h_2)$, the triangle inequality bounds the
norm of each component of $Mh$ by the corresponding component of the scalar
matrix of entry norms applied to $(\norm{h_1},\norm{h_2})$.  The fifth statement follows from
\[
 \norm A_2=\sup_{\norm{x}_2=\norm{y}_2=1}\abs{y^{\mathsf T}Ax}
\]
and the observation that, for $A\geq0$ entrywise,
\(
 \abs{y^{\mathsf T}Ax}\leq |y|^{\mathsf T}A|x|
\), where $|x|$ and $|y|$ denote coordinatewise absolute values.
\end{proof}

Because $H\cong H\oplus H$, there exist isometries $u,v\in\mathcal B(H)$ with
orthogonal ranges whose range projections sum to the identity:
\begin{equation}\label{eq:cuntz}
 u^*u=v^*v=\one,
 \qquad u^*v=v^*u=0,
 \qquad uu^*+vv^*=\one.
\end{equation}
For example, choose a unitary $W:H\oplus H\to H$ and put
$u\xi=W(\xi,0)$ and $v\xi=W(0,\xi)$.  This construction works equally in the
nonseparable case.

\begin{lemma}\label{lem:cuntz-block}
The map $U:H\oplus H\to H$ given by
$U(\xi,\eta)=u\xi+v\eta$ is unitary.  Consequently,
\begin{align}
&\norm{u z_{11}u^*+u z_{12}v^*+v z_{21}u^*+v z_{22}v^*}
\notag\\
&\hspace{45mm}=
\norm{
\begin{pmatrix}
z_{11}&z_{12}\\z_{21}&z_{22}
\end{pmatrix}}_{M_2(\mathcal B(H))}.
\label{eq:block-identification}
\end{align}
In particular,
\begin{equation}\label{eq:row-bound}
 \norm{xv^*+yu^*}
 \leq\bigl(\norm x^2+\norm y^2\bigr)^{1/2}.
\end{equation}
\end{lemma}

\begin{proof}
The relations \eqref{eq:cuntz} show that $U$ is an onto isometry and that the
operator on the left of \eqref{eq:block-identification} is
$U\bigl(\begin{smallmatrix}z_{11}&z_{12}\\z_{21}&z_{22}\end{smallmatrix}\bigr)U^*$.
For \eqref{eq:row-bound}, use
\[
 (xv^*+yu^*)(xv^*+yu^*)^*=xx^*+yy^*.
\]
\end{proof}

\section{The commutator template and diagonal similarity}

Fix an integer $n\geq3$, a number $\delta>0$, and operators
$\beta_1,\ldots,\beta_n\in\mathcal B(H)$.  Set
$\beta_0=\beta_{n+1}=0$.  Define $X,D\in M_n(\mathcal B(H))$ by
\begin{align}
 X_{ij}
 &=\one\ind{i=j+1}+\beta_i\ind{j=n},
 \label{eq:X-def}\\
 D_{ij}
 &=\delta^{-1}u\ind{i=j+1}
 +\delta^{-1}v\ind{i=j}
 +i\one\ind{j=i+1}
 +\delta^{-1}\beta_i u\ind{j=n}.
 \label{eq:D-def}
\end{align}
When two conditions in \eqref{eq:D-def} hold simultaneously, the corresponding
terms are added.

\begin{lemma}[Exact commutator identity]\label{lem:commutator}
For $1\leq i,j\leq n$,
\begin{align}
 [D,X]_{ij}
 &=\one\ind{i=j}
 \notag\\
 &\quad+
 \left[
 \delta^{-1}\bigl([v,\beta_i]+[u,\beta_{i-1}]
 +\beta_i[u,\beta_n]\bigr)
 +i\beta_{i+1}-n\one\ind{i=n}
 \right]\ind{j=n}.
 \label{eq:commutator}
\end{align}
Thus $[D,X]-\one$ is supported in the last column.
\end{lemma}

\begin{proof}
Let $S_{ij}=\one\ind{i=j+1}$, let
$B_{ij}=\beta_i\ind{j=n}$, and let
$N_{ij}=i\one\ind{j=i+1}$.  Then $X=S+B$ and
\begin{equation}\label{eq:D-compact}
 D=\delta^{-1}\diag(v)+\delta^{-1}X\diag(u)+N.
\end{equation}
The commutator $[\diag(v),X]$ vanishes off the last column, where its row-$i$
entry is $[v,\beta_i]$.  Moreover,
\[
 [X\diag(u),X]=X[\diag(u),X],
\]
which also vanishes off the last column, where its row-$i$ entry is
$[u,\beta_{i-1}]+\beta_i[u,\beta_n]$.

Finally, direct multiplication gives
\[
 [N,X]_{ij}
 =\one\ind{i=j}
 +\bigl(i\beta_{i+1}-n\one\ind{i=n}\bigr)\ind{j=n}.
\]
Adding the three contributions in \eqref{eq:D-compact} proves
\eqref{eq:commutator}.
\end{proof}

We say that $\beta=(\beta_1,\ldots,\beta_n)$ solves the nonlinear system if
\begin{equation}\label{eq:system}
 [v,\beta_i]+[u,\beta_{i-1}]
 +\beta_i[u,\beta_n]+\delta i\beta_{i+1}
 =\delta n\one\ind{i=n},
 \qquad 2\leq i\leq n.
\end{equation}
If \eqref{eq:system} holds, then $[D,X]-\one$ has only one
possibly nonzero entry, namely
\begin{equation}\label{eq:single-corner}
 Z_{1n}=\delta^{-1}\bigl([v,\beta_1]+\beta_1[u,\beta_n]\bigr)+\beta_2.
\end{equation}

For $0<\mu\leq1$, put
\begin{equation}\label{eq:Smu}
 \begin{aligned}
 S_\mu&=\diag(\mu^{n-1}\one,\mu^{n-2}\one,\ldots,\one),\\
 D_\mu&=\mu^{-1}S_\mu D S_\mu^{-1},
 &\qquad
 X_\mu&=\mu S_\mu X S_\mu^{-1}.
 \end{aligned}
\end{equation}

\begin{proposition}[Diagonal similarity]\label{prop:similarity}
For arbitrary $\beta_1,\ldots,\beta_n$,
\begin{equation}\label{eq:similarity-comm}
 [D_\mu,X_\mu]=S_\mu[D,X]S_\mu^{-1},
\end{equation}
\begin{equation}\label{eq:Xmu-bound}
 \norm{X_\mu}
 \leq1+\sum_{i=1}^n\mu^{n-i+1}\norm{\beta_i},
\end{equation}
and
\begin{equation}\label{eq:Dmu-bound}
 \norm{D_\mu}
 \leq\frac1{\mu\delta}+\frac1{\mu^2\delta}+(n-1)
 +\frac1\delta\sum_{i=1}^n\mu^{n-i-1}\norm{\beta_i}.
\end{equation}
If $\beta$ solves \eqref{eq:system}, then
\begin{equation}\label{eq:error-conjugated}
 \norm{[D_\mu,X_\mu]-\one}
 =\mu^{n-1}\norm{Z_{1n}}.
\end{equation}
\end{proposition}

\begin{proof}
The scalar factors in \eqref{eq:Smu} cancel in the commutator, which gives
\eqref{eq:similarity-comm}.  Conjugation by $S_\mu$ multiplies an $(i,j)$
entry by $\mu^{j-i}$.  Hence $X_\mu$ is the unweighted subdiagonal shift plus
a last column with row-$i$ entry $\mu^{n-i+1}\beta_i$, proving
\eqref{eq:Xmu-bound} by Lemma~\ref{lem:matrix-norms}.

Similarly, the four parts of $D_\mu$ are a diagonal of norm
$(\mu\delta)^{-1}$, a subdiagonal of norm $(\mu^2\delta)^{-1}$, a weighted
superdiagonal of norm $n-1$, and a last column with row-$i$ entry
$\delta^{-1}\mu^{n-i-1}\beta_i u$.  This gives
\eqref{eq:Dmu-bound}.  If \eqref{eq:system} holds, then
\eqref{eq:single-corner} and the first part of Lemma~\ref{lem:matrix-norms} give
\eqref{eq:error-conjugated}.
\end{proof}

\section{A right inverse for the linearized system}

Let
\[
 Y:=\mathcal B(H)^{n-1},\qquad
 Z:=\mathcal B(H)^n,
\]
where elements of $Y$ are indexed by $2,\ldots,n$, and both spaces carry the
maximum norm.  For $x\in Y$ we use the boundary conventions
$x_1=x_{n+1}=0$.  Define $T:Z\to Y$, $L:Y\to Z$, and $E:Y\to Y$ by
\begin{align}
 (Tb)_i&=[v,b_i]+[u,b_{i-1}],
 &&2\leq i\leq n,
 \label{eq:T-def}\\
 (Lx)_i&=-\frac12\bigl(x_i v^*+x_{i+1}u^*\bigr),
 &&1\leq i\leq n,
 \label{eq:L-def}\\
 (Ex)_i&=\frac12\bigl(vx_iv^*+vx_{i+1}u^*
 +ux_{i-1}v^*+ux_iu^*\bigr),
 &&2\leq i\leq n.
 \label{eq:E-def}
\end{align}

\begin{lemma}\label{lem:TL}
One has
\begin{equation}\label{eq:TL}
 TL=I_Y-E.
\end{equation}
\end{lemma}

\begin{proof}
For $2\leq i\leq n$, the Cuntz relations give
\begin{align*}
 [v,(Lx)_i]
 &=-\frac12(vx_iv^*+vx_{i+1}u^*)+\frac12x_i,\\
 [u,(Lx)_{i-1}]
 &=-\frac12(ux_{i-1}v^*+ux_iu^*)+\frac12x_i.
\end{align*}
Their sum is $x_i-(Ex)_i$.
\end{proof}

\begin{lemma}[Weighted contraction]\label{lem:weighted-contraction}
For $x\in Y$, define
\begin{equation}\label{eq:prime-norm}
 \norm x'
 :=\max_{2\leq i\leq n}
 \left(2-\frac{i^2}{n^2}\right)^{-1/2}\norm{x_i}.
\end{equation}
Then
\begin{equation}\label{eq:norm-equivalence}
 \frac1{\sqrt2}\norm x_Y\leq\norm x'\leq\norm x_Y
\end{equation}
and
\begin{equation}\label{eq:E-contraction}
 \norm{Ex}'
 \leq\left(1-\frac1{8n^2}\right)\norm x'.
\end{equation}
\end{lemma}

\begin{proof}
The norm equivalence is immediate from
$1\leq2-i^2/n^2\leq2$ for $2\leq i\leq n$.
By homogeneity, assume $\norm x'\leq1$, so
$\norm{x_i}^2\leq2-i^2/n^2$.  This inequality also holds for $i=1,n+1$
because $x_1=x_{n+1}=0$; note that
$2-(n+1)^2/n^2>0$ for $n\geq3$.

By Lemma~\ref{lem:cuntz-block},
\begin{equation}\label{eq:E-block}
 \norm{(Ex)_i}
 =\frac12\norm{
 \begin{pmatrix}
 x_i&x_{i-1}\\x_{i+1}&x_i
 \end{pmatrix}}_{M_2(\mathcal B(H))}.
\end{equation}
Using Lemma~\ref{lem:matrix-norms} and then the Frobenius norm of the scalar matrix
of entry norms,
\begin{align*}
 \norm{(Ex)_i}^2
 &\leq\frac14\bigl(2\norm{x_i}^2+
 \norm{x_{i-1}}^2+\norm{x_{i+1}}^2\bigr)\\
 &\leq 2-\frac{i^2}{n^2}-\frac1{2n^2}\\
 &\leq\left(1-\frac1{4n^2}\right)
 \left(2-\frac{i^2}{n^2}\right)\\
 &\leq\left(1-\frac1{8n^2}\right)^2
 \left(2-\frac{i^2}{n^2}\right).
\end{align*}
Taking square roots and maximizing over $i$ proves
\eqref{eq:E-contraction}.
\end{proof}

\begin{proposition}[Right inverse]\label{prop:right-inverse}
The operator $I_Y-E$ is invertible, and
\begin{equation}\label{eq:R-def}
 R:=L(I_Y-E)^{-1}:Y\to Z
\end{equation}
is a right inverse for $T$.  Moreover,
\begin{equation}\label{eq:R-global}
 \norm R_{Y\to Z}\leq8\sqrt2\,n^2.
\end{equation}
\end{proposition}

\begin{proof}
The Neumann series converges in the primed norm by
Lemma~\ref{lem:weighted-contraction}, and
\[
 \norm{(I_Y-E)^{-1}}'\leq8n^2.
\]
Thus Lemma~\ref{lem:TL} gives $TR=I_Y$.  If $\norm x_Y\leq1$ and
$y=(I_Y-E)^{-1}x$, then
$\norm{y_i}\leq8\sqrt2 n^2$.  By \eqref{eq:L-def} and
\eqref{eq:row-bound},
\[
 \norm{(Rx)_i}
 \leq\frac12\bigl(\norm{y_i}^2+\norm{y_{i+1}}^2\bigr)^{1/2}
 \leq8n^2,
\]
which is stronger than \eqref{eq:R-global}.
\end{proof}

\begin{remark}
The $n^2$ scale in \eqref{eq:R-global} agrees with the spectral-gap scale of
the corresponding scalar Dirichlet averaging operator.  We do not need a
matching lower bound for $R$; the improvement below comes from estimating the
specific structured data fed into $R$, rather than its global norm.
\end{remark}

\section{A monotone supersolution principle}

For a nonnegative vector $p=(p_1,\ldots,p_{n+1})$, define
\begin{equation}\label{eq:A-def}
 (Ap)_i
 :=\frac12\norm{
 \begin{pmatrix}
 p_i&p_{i-1}\\p_{i+1}&p_i
 \end{pmatrix}}_{M_2(\mathbb R)},
 \qquad 2\leq i\leq n.
\end{equation}
By the last part of Lemma~\ref{lem:matrix-norms}, $A$ is monotone in every
coordinate.

\begin{lemma}[Majorant property]\label{lem:majorant}
Let $x\in Y$ and put $p_j=\norm{x_j}$, with $p_1=p_{n+1}=0$.  Then
\[
 \norm{(Ex)_i}\leq(Ap)_i,
 \qquad 2\leq i\leq n.
\]
\end{lemma}

\begin{proof}
Combine \eqref{eq:E-block} with the fourth part of
Lemma~\ref{lem:matrix-norms}.
\end{proof}

The following elementary inequality is the key quantitative input.

\begin{lemma}[Sharp two-by-two estimate]\label{lem:two-by-two}
Let $a,b,c\geq0$ and assume $b^2+c^2\leq2a^2$.  Then
\begin{equation}\label{eq:two-by-two}
 \norm{
 \begin{pmatrix}a&b\\c&a\end{pmatrix}}_{M_2(\mathbb R)}
 \leq a+\sqrt{\frac{b^2+c^2}{2}}.
\end{equation}
Consequently, if $a>0$ and
$d:=a^2-(b^2+c^2)/2\geq0$, then
\begin{equation}\label{eq:two-by-two-defect}
 \frac12\norm{
 \begin{pmatrix}a&b\\c&a\end{pmatrix}}
 \leq a-\frac d{4a}.
\end{equation}
Equality in \eqref{eq:two-by-two} is possible, for example when $b=c$.
\end{lemma}

\begin{proof}
Put $P=a^2$, $Q=b^2$, and $R=c^2$.  A direct computation of the largest
eigenvalue of $M^{\mathsf T}M$, where
$M=\bigl(\begin{smallmatrix}a&b\\c&a\end{smallmatrix}\bigr)$, gives
\begin{equation}\label{eq:Mnorm-exact}
 \norm M^2
 =P+\frac{Q+R}{2}
 +\sqrt{\frac{(Q-R)^2}{4}+P(b+c)^2}.
\end{equation}
Let $W=(b+c)^2$ and $e=(b-c)^2$.  Then
$W+e=2(Q+R)$ and $(Q-R)^2=eW$.  Since $Q+R\leq2P$, one has
$W\leq4P$, and therefore
\[
 \frac{eW}{4}+PW\leq eP+PW=P(e+W)=2P(Q+R).
\]
Substitution into \eqref{eq:Mnorm-exact} gives
\[
 \norm M^2
 \leq P+\frac{Q+R}{2}+\sqrt{2P(Q+R)}
 =\left(\sqrt P+\sqrt{\frac{Q+R}{2}}\right)^2,
\]
which proves \eqref{eq:two-by-two}.

For \eqref{eq:two-by-two-defect}, write
$(Q+R)/2=P-d$ and use the tangent-line bound
$\sqrt{P-d}\leq\sqrt P-d/(2\sqrt P)$.
\end{proof}

\begin{proposition}[Supersolution bound]\label{prop:supersolution}
Let $s=(s_i)_{i=2}^n$ be nonnegative and let
$\psi=(\psi_1,\ldots,\psi_{n+1})$ be nonnegative.  Suppose
\begin{equation}\label{eq:supersolution}
 s_i+(A\psi)_i\leq\psi_i,
 \qquad 2\leq i\leq n.
\end{equation}
If $x\in Y$ satisfies $\norm{x_i}\leq s_i$, then
\begin{equation}\label{eq:neumann-profile}
 \norm{\left(\sum_{k=0}^{K}E^kx\right)_i}\leq\psi_i
 \qquad(K\geq0,
 \ 2\leq i\leq n),
\end{equation}
\begin{equation}\label{eq:resolvent-profile}
 \norm{((I_Y-E)^{-1}x)_i}\leq\psi_i,
\end{equation}
and
\begin{equation}\label{eq:R-profile}
 \norm{(Rx)_i}
 \leq\frac12\bigl(\psi_i^2+\psi_{i+1}^2\bigr)^{1/2},
 \qquad 1\leq i\leq n.
\end{equation}
\end{proposition}

\begin{proof}
Set $y^{[K]}=\sum_{k=0}^K E^kx$.  The case $K=0$ follows from
$s_i\leq\psi_i$.  If \eqref{eq:neumann-profile} holds for $K$, then
$y^{[K+1]}=x+Ey^{[K]}$, and the majorant property together with the monotonicity
of $A$ gives
\[
 \norm{y_i^{[K+1]}}
 \leq s_i+(A\psi)_i\leq\psi_i.
\]
This proves \eqref{eq:neumann-profile} by induction.  Passing to the Neumann
series limit proves \eqref{eq:resolvent-profile}.  Finally, apply
\eqref{eq:row-bound} to $R=L(I_Y-E)^{-1}$.
\end{proof}

\begin{corollary}[Concavity criterion]\label{cor:concavity}
Assume $\psi_i>0$ for $1\leq i\leq n$, put $\psi_{n+1}=0$, and define
$\Psi_i=\psi_i^2$ and
\begin{equation}\label{eq:defect}
 d_i:=\Psi_i-\frac{\Psi_{i-1}+\Psi_{i+1}}2,
 \qquad 2\leq i\leq n.
\end{equation}
If
\begin{equation}\label{eq:concavity-criterion}
 d_i\geq0
 \qquad\text{and}\qquad
 4\psi_i s_i\leq d_i,
 \qquad 2\leq i\leq n,
\end{equation}
then \eqref{eq:supersolution} holds, and all conclusions of
Proposition~\ref{prop:supersolution} follow.
\end{corollary}

\begin{proof}
Apply \eqref{eq:two-by-two-defect} with
$a=\psi_i$, $b=\psi_{i-1}$, and $c=\psi_{i+1}$.  The first condition in
\eqref{eq:concavity-criterion} is exactly
$b^2+c^2\leq2a^2$, and it yields
\[
 (A\psi)_i\leq\psi_i-\frac{d_i}{4\psi_i}
 \leq\psi_i-s_i.
\]
\end{proof}

\section{Solving the nonlinear system}

From now on let
\begin{equation}\label{eq:delta-choice}
 n\geq9,
 \qquad
 \delta:=\frac1{400n^3}.
\end{equation}
Define
\begin{align}
 \Psi_i&:=144\delta^2 i(2n-i),
 &&1\leq i\leq n,
 &\Psi_{n+1}&:=0,
 \label{eq:Psi-profile}\\
 \psi_i&:=\sqrt{\Psi_i}
 =12\delta\sqrt{i(2n-i)},
 &&1\leq i\leq n,
 &\psi_{n+1}&:=0,
 \label{eq:psi-profile}\\
 \overline b_i&:=\frac12\sqrt{\Psi_i+\Psi_{i+1}},
 &&1\leq i\leq n,
 &\overline b_{n+1}&:=0,
 \label{eq:bbar}\\
 s_i&:=\delta n\ind{i=n}
 +2\overline b_i\overline b_n
 +\delta i\overline b_{i+1},
 &&2\leq i\leq n.
 \label{eq:s-def}
\end{align}

\begin{lemma}[Profile estimates]\label{lem:profile-estimates}
The following estimates hold:
\begin{align}
 \psi_i&\leq12\delta n,
 \label{eq:profile-p1}\\
 \psi_i&\leq12\sqrt2\,\delta\sqrt{ni}
 <17\delta\sqrt{ni},
 \label{eq:profile-p2}\\
 \overline b_i&\leq6\sqrt2\,\delta n,
 \label{eq:profile-p3}\\
 \overline b_i&\leq15\delta\sqrt{ni},
 \qquad
 \overline b_n=6\delta n,
 \label{eq:profile-p4}\\
 d_i&=144\delta^2,
 \qquad 2\leq i\leq n-1,
 \label{eq:profile-p5a}\\
 d_n&=72\delta^2(n^2+1),
 \label{eq:profile-p5b}\\
 \overline b_1&\leq6\sqrt6\,\delta\sqrt n,
 \qquad
 \overline b_2\leq6\sqrt{10}\,\delta\sqrt n.
 \label{eq:profile-p6}
\end{align}
\end{lemma}

\begin{proof}
The function $i(2n-i)=n^2-(n-i)^2$ is increasing on $[0,n]$ and bounded by
$n^2$, giving \eqref{eq:profile-p1}.  Since
$i(2n-i)\leq2ni$, one obtains \eqref{eq:profile-p2}.

From $4\overline b_i^2=\Psi_i+\Psi_{i+1}\leq288\delta^2n^2$ one obtains
\eqref{eq:profile-p3}.  For $i<n$,
\[
 4\overline b_i^2
 \leq144\delta^2\bigl(2ni+2n(i+1)\bigr)
 \leq864\delta^2ni,
\]
which implies $\overline b_i\leq\sqrt{216}\,\delta\sqrt{ni}<15\delta\sqrt{ni}$.
The formula for $\overline b_n$ follows from $\Psi_{n+1}=0$.

For $2\leq i<n$, the second difference of $i(2n-i)$ is $-2$, hence
$d_i=144\delta^2$.  At $i=n$,
\[
 d_n
 =144\delta^2n^2-\frac12\,144\delta^2(n^2-1)
 =72\delta^2(n^2+1).
\]
Finally,
\[
 \Psi_1+\Psi_2=144\delta^2(6n-5)
 \leq864\delta^2n
\]
and
\[
 \Psi_2+\Psi_3=144\delta^2(10n-13)
 \leq1440\delta^2n,
\]
which give \eqref{eq:profile-p6}.
\end{proof}

\begin{lemma}[Verification of the supersolution condition]\label{lem:profile-super}
For every $2\leq i\leq n$,
\begin{equation}\label{eq:profile-super}
 4\psi_i s_i\leq d_i.
\end{equation}
Consequently, the pair $(s,\psi)$ satisfies
\eqref{eq:supersolution}.
\end{lemma}

\begin{proof}
For $i=n$, \eqref{eq:s-def} and \eqref{eq:profile-p4} give
\[
 s_n=\delta n+72\delta^2n^2,
 \qquad
 \psi_n=12\delta n.
\]
Therefore
\[
 4\psi_ns_n
 =48\delta^2n^2+3456\delta^3n^3
 =\left(48n^2+\frac{3456}{400}\right)\delta^2
 \leq72(n^2+1)\delta^2=d_n.
\]

Now let $2\leq i\leq n-1$.  By
\eqref{eq:profile-p4},
\begin{align*}
 s_i
 &\leq2\bigl(15\delta\sqrt{ni}\bigr)(6\delta n)
 +\delta i\bigl(15\sqrt2\,\delta\sqrt{ni}\bigr)\\
 &\leq201.3\,\delta^2n\sqrt{ni}.
\end{align*}
Together with \eqref{eq:profile-p2}, this gives
\[
 4\psi_i s_i
 \leq13689\,\delta^3n^2i
 \leq13689\,\delta^3n^3
 =\frac{13689}{400}\delta^2
 <144\delta^2=d_i.
\]
The conclusion follows from Corollary~\ref{cor:concavity}.
\end{proof}

\begin{proposition}[Nonlinear solution]\label{prop:nonlinear-solution}
For every $n\geq9$ and $\delta=1/(400n^3)$, the system
\eqref{eq:system} has a solution
$\beta=(\beta_1,\ldots,\beta_n)\in\mathcal B(H)^n$ such that
\begin{equation}\label{eq:beta-bounds}
 \norm{\beta_i}\leq\overline b_i,
 \qquad 1\leq i\leq n.
\end{equation}
In particular,
\begin{equation}\label{eq:beta-consequences}
 \begin{aligned}
 \max_i\norm{\beta_i}&\leq6\sqrt2\,\delta n,
 &\qquad \norm{\beta_n}&\leq6\delta n,\\
 \norm{\beta_1}&\leq6\sqrt6\,\delta\sqrt n,
 &\qquad \norm{\beta_2}&\leq6\sqrt{10}\,\delta\sqrt n.
 \end{aligned}
\end{equation}
\end{proposition}

\begin{proof}
Let
\[
 Q:=\{y\in Y:\norm{y_i}\leq s_i\text{ for }2\leq i\leq n\}.
\]
This is a nonempty closed subset of the Banach space $Y$, hence complete.
Define $\Phi:Y\to Y$ by
\begin{equation}\label{eq:Phi}
 \Phi(y)_i
 :=\delta n\one\ind{i=n}
 -(Ry)_i[u,(Ry)_n]
 -\delta i(Ry)_{i+1},
 \qquad 2\leq i\leq n,
\end{equation}
where $(Ry)_{n+1}=0$.

If $y\in Q$ and $b=Ry$, then
Lemma~\ref{lem:profile-super} and Proposition~\ref{prop:supersolution} give
\[
 \norm{b_i}
 \leq\frac12\sqrt{\psi_i^2+\psi_{i+1}^2}
 =\overline b_i.
\]
Hence
\[
 \norm{\Phi(y)_i}
 \leq\delta n\ind{i=n}
 +2\overline b_i\overline b_n
 +\delta i\overline b_{i+1}
 =s_i,
\]
so $\Phi(Q)\subseteq Q$.

For $y,y'\in Q$, put $b=Ry$ and $b'=Ry'$.  Using the preceding
profile estimate together with \eqref{eq:profile-p3} and
$\overline b_n=6\delta n$, one obtains
\begin{align*}
 \norm{\Phi(y)-\Phi(y')}_Y
 &\leq(12\delta n+12\sqrt2\,\delta n+\delta n)
 \norm{b-b'}_Z\\
 &\leq(13+12\sqrt2)\delta n\,
 8\sqrt2\,n^2\norm{y-y'}_Y\\
 &=\frac{192+104\sqrt2}{400}\norm{y-y'}_Y
 <0.85\norm{y-y'}_Y.
\end{align*}
Thus $\Phi$ is a contraction of $Q$.  Let $y^*\in Q$ be its unique fixed
point and put $\beta=Ry^*$.  Since $TR=I_Y$,
\[
 T\beta=y^*=\Phi(y^*),
\]
which is exactly \eqref{eq:system}.  The bounds follow from the profile
estimate above and Lemma~\ref{lem:profile-estimates}.
\end{proof}

\section{Norm estimates and completion of the proof}

Take $\mu=1/2$ in Proposition~\ref{prop:similarity}, and let $D_\mu,X_\mu$ be built from
the solution in Proposition~\ref{prop:nonlinear-solution}.

\begin{proposition}\label{prop:final-estimates}
For $n\geq9$ and $\delta=1/(400n^3)$,
\begin{align}
 \norm{X_\mu}&\leq1.00027,
 \label{eq:final-X}\\
 \norm{D_\mu}&\leq2401n^3,
 \label{eq:final-D}\\
 \norm{D_\mu}\norm{X_\mu}&\leq2402n^3,
 \label{eq:final-product}\\
 \norm{[D_\mu,X_\mu]-\one}&\leq59\sqrt n\,2^{-n}.
 \label{eq:final-error}
\end{align}
\end{proposition}

\begin{proof}
Let $\overline b=\max_i\overline b_i$.  By
\eqref{eq:profile-p3},
$\overline b\leq6\sqrt2\,\delta n$.  Since
$\sum_{i=1}^n2^{-(n-i+1)}<1$, \eqref{eq:Xmu-bound} gives
\[
 \norm{X_\mu}
 \leq1+\overline b
 \leq1+\frac{6\sqrt2}{400n^2}
 \leq1.00027.
\]

For $D_\mu$, \eqref{eq:Dmu-bound} and
$\sum_{i=1}^n2^{i+1-n}<4$ yield
\begin{align*}
 \norm{D_\mu}
 &\leq\frac6\delta+(n-1)+\frac{4\overline b}{\delta}\\
 &\leq2400n^3+(1+24\sqrt2)n\\
 &<2400n^3+35n
 \leq2401n^3.
\end{align*}
The product estimate follows because
$2401\cdot1.00027<2402$.

It remains to estimate the surviving corner in
\eqref{eq:single-corner}.  By \eqref{eq:beta-consequences},
\begin{align*}
 \norm{Z_{1n}}
 &\leq\frac2\delta\norm{\beta_1}
 +\frac2\delta\norm{\beta_1}\norm{\beta_n}
 +\norm{\beta_2}\\
 &\leq\sqrt n\left(
 12\sqrt6+72\sqrt6\,\delta n+6\sqrt{10}\,\delta
 \right).
\end{align*}
For $n\geq9$, one has
$\delta n\leq1/32400$ and $\delta\leq1/291600$.  Using
$\sqrt6<2.4495$ and $\sqrt{10}<3.1623$, the coefficient in parentheses is
strictly less than $29.4$.  Hence, by \eqref{eq:error-conjugated},
\[
 \norm{[D_\mu,X_\mu]-\one}
 \leq29.4\sqrt n\,2^{-(n-1)}
 <59\sqrt n\,2^{-n}.
\]
\end{proof}

\begin{proof}[Proof of Theorem~\ref{thm:main}]
Let
\[
 t:=\log_2\frac1\varepsilon\geq1,
 \qquad
 n:=\max\left\{9,\,\lceil2t\rceil+6\right\}.
\]
Construct $D_\mu,X_\mu\in M_n(\mathcal B(H))$ as above and transport them to
$\mathcal B(H)$ through any unitary $H^{\oplus n}\to H$.  This preserves
norms, products, the identity, and commutators.

We first prove the error bound.  If $n=9$, then
$\lceil2t\rceil\leq3$, so $t\leq3/2$, and
\[
 59\sqrt n\,2^{-n}
 =\frac{177}{512}
 <2^{-3/2}
 \leq2^{-t}=\varepsilon.
\]
If $n\geq10$, then $n=\lceil2t\rceil+6$, so
$t\leq(n-6)/2$.  Therefore
\[
 \frac{59\sqrt n\,2^{-n}}{\varepsilon}
 =59\sqrt n\,2^{t-n}
 \leq\frac{59\sqrt n}{8\,2^{n/2}}
 \leq\frac{59\sqrt{10}}{256}<1,
\]
where the middle expression decreases for integers $n\geq10$.
Thus \eqref{eq:main-error} follows from \eqref{eq:final-error}.

For the norm product, note that $n\leq9t$: this is immediate when $n=9$, and
otherwise
$n\leq2t+7\leq9t$ because $t\geq1$.  Writing
$L=\log(1/\varepsilon)=t\log2$, \eqref{eq:final-product} gives
\begin{align*}
 \norm{D_\varepsilon}\norm{X_\varepsilon}
 &\leq2402n^3\\
 &\leq2402\left(\frac9{\log2}\right)^3L^3\\
 &<5.3\times10^6L^3.
\end{align*}
This proves \eqref{eq:main-norm} and completes the proof.
\end{proof}

\section{Further remarks}

\begin{remark}[Finite-dimensional obstruction]
If $\dim H=N<\infty$, then $\operatorname{tr}[D,X]=0$, so
\[
 \norm{[D,X]-\one}
 \geq\frac1N\abs{\operatorname{tr}([D,X]-\one)}=1.
\]
Thus the approximation problem is genuinely infinite-dimensional.  The proof
uses infinite dimensionality precisely through the identifications
$H\cong H^{\oplus n}$ and the Cuntz relations \eqref{eq:cuntz}.
\end{remark}

\begin{remark}[Effectivity]
The solution in Proposition~\ref{prop:nonlinear-solution} is obtained by a Neumann series
and a Picard iteration with an explicit contraction factor below $0.85$.
Truncating both procedures therefore gives finite noncommutative polynomials in
$u,v,u^*,v^*$ with explicitly controlled residuals.  In particular, the
construction is algorithmic; the limiting formulation is used only to keep the
main proof concise.
\end{remark}

\begin{remark}[Source of the exponent]
The profile estimate gives $\norm{\beta_i}=O(\delta n)$.  The quadratic and
shift terms in \eqref{eq:system} are both of size $O(\delta^2n^2)$.  A profile
of height $O(\delta n)$ over a chain of length $n$ has an interior concavity
budget of order $\delta^2$, so absorption requires
$\delta^3n^3\lesssim\delta^2$, equivalently $\delta\lesssim n^{-3}$.  Since
$\norm D\norm X=O(\delta^{-1})$ and
$n\asymp\log(1/\varepsilon)$, this produces the cubic logarithm.
\end{remark}

\begin{remark}[What remains open]
Combining \eqref{eq:popa} with Theorem~\ref{thm:main} leaves
\[
 \frac12\log\frac1\varepsilon
 \leq m_H(\varepsilon)
 \leq C\left(\log\frac1\varepsilon\right)^3.
\]
No lower-bound improvement is claimed here.  Reaching the conjectural linear
logarithmic scale would require a further structural gain beyond the present
endpoint supersolution estimate.
\end{remark}

\begin{thebibliography}{10}

\bibitem{Bilich}
B. Bilich,
\emph{An $O(\log^4(1/\varepsilon))$ refinement of Tao's construction of
commutators close to the identity},
unpublished manuscript and Lean formalization, 2026,
\url{https://github.com/bilichboris/TaoCommutators}.

\bibitem{BrownPearcy}
A. Brown and C. Pearcy,
\emph{Structure of commutators of operators},
Ann. of Math. (2) \textbf{82} (1965), 112--127.

\bibitem{KrishnaJohnson}
K. M. Krishna and P. S. Johnson,
\emph{Commutators close to the identity in unital $C^*$-algebras},
Proc. Indian Acad. Sci. Math. Sci. \textbf{132} (2022), Paper No. 11.

\bibitem{Popa}
S. Popa,
\emph{On commutators in properly infinite $W^*$-algebras},
in \emph{Invariant Subspaces and Other Topics},
Operator Theory: Advances and Applications, vol. 6,
Birkh\"auser, Basel, 1982, pp. 195--207.

\bibitem{Tao}
T. Tao,
\emph{Commutators close to the identity},
J. Operator Theory \textbf{82} (2019), no. 2, 369--382.

\bibitem{Wielandt}
H. Wielandt,
\emph{\"Uber die Unbeschr\"anktheit der Operatoren der Quantenmechanik},
Math. Ann. \textbf{121} (1949/50), 21.

\bibitem{Wintner}
A. Wintner,
\emph{The unboundedness of quantum-mechanical matrices},
Phys. Rev. (2) \textbf{71} (1947), 738--739.

\end{thebibliography}

\end{document}
